% ==========================================================================
% GENERATED by scripts/split_paper.py. Do not edit this file.
%
% Source: imported_papers/Koopman-JEPA-TMLR-rewrite/main.tex
% Part 12: Why Closed-Form Control Does Not Follow
%
% Edits belong upstream, in the published paper. To change what the
% chapter includes, change imports/bounded.toml and re-run
% the splitter.
% ==========================================================================
\section{Why Closed-Form Control Does Not Follow}
\label{app:lqr}
% [orig app + orig 6 + draft 4.4 + new] DONE. The proposition moves here from
% the approved paper's Section 6, which is now the second paragraph of Section
% 5.5, so the statement and its proof are together for the first time.
This appendix checks the claim of Section~\ref{sec:degeneracy} that ruling out divergent modes does not deliver closed-form control.

\begin{proposition}[No divergent eigenmodes under orthogonal $\K$]
\label{prop:cayley-stab}
Let $\K \in \mathbb{R}^{D\times D}$ be the Cayley-parameterized transition matrix defined by Equation~\ref{eq:cayley} for any learnable weight matrix $W \in \mathbb{R}^{D\times D}$. Then every eigenvalue satisfies $|\lambda_i| = 1$, so $\K$ has no divergent mode. The DARE failure mode in which a divergent mode lies outside the subspace reachable through $\B$ is therefore eliminated for every choice of action matrix $\B$, and stabilizability of $(\K, \B)$ reduces to the single remaining condition that the unit-circle modes are controllable through $\B$~\citep{Lancaster1995}. This is a property of the parameterization rather than of the learned weights, and it holds for any $W$.
\end{proposition}

\paragraph{Goal, coordinate shift, and error dynamics.}
We take $\z^*$ to be a fixed point of the latent dynamics under zero action up to an offset $d = \K\z^* - \z^*$, which enters the error dynamics as a constant disturbance. For the Cayley variant this offset is structural, since $\K\z = \z$ holds exactly only on $\ker A$, which is trivial for a generic skew-symmetric $A$ in even dimension. Because standard LQR regulates the system toward the origin, we apply a coordinate shift and define the error state $\tilde{\z}_t := \z_t - \z^*$. This is the quantity LQR drives to zero, and its infinite-horizon cost is the quadratic functional $\sum_{t=0}^{\infty}(\tilde{\z}_t^\top Q \tilde{\z}_t + \mathbf{a}_t^\top R \mathbf{a}_t)$, so minimizing it pushes the system back toward $\z^*$ while penalizing actuation effort. The error dynamics are approximated as
\begin{equation}
  \tilde{\z}_{t+1} \approx \K\tilde{\z}_t + \B\mathbf{a}_t,
  \label{eq:error_dynamics}
\end{equation}
where $\B \in \mathbb{R}^{D \times m}$ is the learned action matrix of Section~\ref{sec:variants}. The cost is minimized by solving the discrete algebraic Riccati equation for a unique positive-definite $P \in \mathbb{R}^{D \times D}$:
\begin{equation}
P = Q + \K^\top P \K - \K^\top P \B (R + \B^\top P \B)^{-1} \B^\top P \K.
\label{eq:dare}
\end{equation}
$Q$ is the state penalty and $R$ the control penalty, both user-defined, and we use $Q = I_D$ and $R = 0.1\,I_m$ throughout. From $P$ the optimal feedback gain is $\mathbf{G} = (R + \B^\top P \B)^{-1} \B^\top P \K$, and the control law, clipped to the permissible action range, is
\begin{equation}
  \mathbf{a}_t = \mathrm{clip}(-\mathbf{G}(\z_t - \z^*),\; -1,\; 1).
  \label{eq:lqr_law}
\end{equation}

\begin{proof}[Proof of Proposition~\ref{prop:cayley-stab}]
By construction $\K^\top \K = I$ (Appendix~\ref{app:dmd}), so $|\lambda_i(\K)| = 1$ for every eigenvalue. No mode therefore lies strictly outside the unit circle, eliminating the source of unbounded $\|\mathbf{G}\|_F$ that arises when divergent eigenvectors fall outside the column space reachable through $\B$. The remaining sufficient conditions for a unique positive-definite DARE solution then reduce to controllability of the unit-circle modes through $\B$, together with positive-definiteness of $(Q, R)$~\citep{Lancaster1995}.
\end{proof}

\paragraph{Solving the Riccati equation so that it can fail.}
Equation~\ref{eq:dare} has no stabilizing solution when the pair $(\K, \B)$ is not stabilizable, and a numerical method has to be able to say so. Iterating the Riccati recursion for a fixed number of steps cannot: it diverges in exactly that case and still returns whatever matrix it has reached, from which a finite-looking gain can be read. We therefore solve with a method that reports failure, report $\|\mathbf{G}\|_F$ only for the seeds where a stabilizing solution exists, and report the closed-loop spectral radius $\rho(\K - \B\mathbf{G})$ in every case, since that is the quantity which says whether the controller stabilizes the learned dynamics at all.

\input{tables/table_control.tex}

\paragraph{What the diagnostics say.}
Table~\ref{tab:control} reports the result. The unconstrained operator admits no stabilizing solution on any of the ten seeds in either environment, which is the failure Proposition~\ref{prop:cayley-stab} identifies and the Cayley parameterization removes. The Cayley operator does remove it, solving on all ten seeds in both environments. What it delivers is a closed loop whose spectral radius is $0.99992$ on pendulum and $0.99996$ on cartpole, so a disturbance halves in about $8{,}600$ and $19{,}000$ steps against rollouts of $50$. The soft variant is the only one that contracts usefully, halving in $51$ steps on pendulum, and it does so on nine seeds of ten while failing to contract at all on cartpole. The reason is visible in the two conditioning columns. With every mode on the unit circle, stabilizability requires all $D{=}256$ modes to be controllable through a scalar input, and the smallest PBH margin is of order $10^{-5}$ throughout while the Krylov matrix is numerically singular in double precision. We report the condition number as a bound rather than a value, because at that magnitude its digits are not meaningful and do not reproduce between runs on identical checkpoints.

% =============================================================================
